\documentclass[10pt,a4paper]{amsart}
\usepackage[margin=19mm]{geometry}
\usepackage{amsmath,amssymb,mathtools}
\usepackage[scr=rsfs,cal=euler]{mathalfa}
\usepackage{xfrac}
\usepackage[unicode,hidelinks,bookmarks=false]{hyperref}
\newcommand{\ie}{i.e.\ }
\newcommand{\cf}{cf.\ }
\newcommand{\resp}{respectively\ }
\newcommand{\Subsection}{\subsection}
\providecommand{\nobreakdash}{\nobreak\hbox{-}}
\setlength{\emergencystretch}{2em}
\multlinegap0pt
\usepackage{harpoon}
\usepackage{tikz}
\usepackage{float}
\usepackage{mathrsfs}
\usepackage{verbatim}
\usepackage{enumitem}
\usetikzlibrary{positioning,chains,fit,shapes,calc}  
\newtheorem{theorem}{Theorem}
\newtheorem{lemma}{Lemma}
\newtheorem{definition}{Definition}
\newtheorem{construction}{Construction}
\newtheorem{remark}{Remark}
\newtheorem{corollary}{Corollary}
\newtheorem{property}{Property}
\newtheorem{requirement}{Requirement}
\newcommand{\vect}[1]{\accentset{\rightharpoonup}{#1}}
\newcommand{\harpoon}{\overset{\rightharpoonup}}
\newcommand{\ppp}{\[\begin{aligned}}
\newcommand{\ooo}{\end{aligned}\]}
\newcommand{\err}{\mathrm{err}}
\newcommand{\rad}{\mathrm{rad}}
\newcommand{\val}{\mathrm{val}}
\newcommand{\ind}{\mathrm{ind}}
\newcommand{\mrm}[1]{\mathrm{#1}}
\newcommand{\sgn}{\mathrm{sgn}}
\newcommand{\degu}{\mathrm{deg}_{u_3}}
\begin{document}
\title[On Boolean polynomials and the Union-Closed Conjecture]{On Boolean polynomials and the Union-Closed Conjecture}
\author{Mario DeFranco}
\date{}
\maketitle
\noindent\emph{Source and licence.} On Boolean polynomials and the Union-Closed Conjecture. Version arXiv:2606.26191v1 (CC BY 4.0). This material is distributed under the Creative Commons Attribution 4.0 International licence, http://creativecommons.org/licenses/by/4.0/, as recorded in the arXiv OAI licence metadata for this exact version and shown on the abstract page https://arxiv.org/abs/2606.26191v1. Full rights notice: licenses/arxiv-2606.26191-CC-BY-4.0.txt.\par\smallskip
\noindent\textit{Editorial scope.} Counted unit: the original \texttt{theorem} environment \texttt{-} at source lines 181--183 together with its complete proof at lines 184--212; the delivered document keeps source lines 86--214 so that every referenced label and every citation resolves inside it. Native editable source is the paper's own concatenated TeX; the AMS wrapper is editorial formatting only. No publisher class, upstream script or unrelated artwork is redistributed.
\section{The Boolean polynomial $\mathrm{ICC}_{m,n}(X)$ }
\begin{definition} 
Fix integers $m,n \geq 2$. Consider a set of $mn$ Boolean variables $$\{x_{i,j} : 1 \leq i \leq m; 1 \leq j \leq n \},$$ i.e. $x_{i,j}^2 = x_{i,j}$. Let $X_i$ denote the sequence  
$$X_i = (x_{i,1}, \ldots, x_{i,n}).$$ 
Let $B$ denote the Boolean algebra $$B=\mathbb{F}_2[\{x_{i,j}\}_{i,j}] .$$


For each specialization $S$ of Boolean variables  to values in $\mathbb{F}_2$, each sequence $X_i$ corresponds to a subset of $\{ 1, \ldots, n\}$, so we will denote such a set by $X_i|_S$. We will also let $u_j, v_j$ for $1 \leq j  \leq n$ denote Boolean variables, and we will denote specializations of $u_j,v_j$ to $\mathbb{F}_2$ by $S$ also. We will call any polynomial of Boolean variables with coefficients in $\mathbb{F}_2$ a \emph{``Boolean polynomial"}. 
\end{definition} 

\begin{definition} For the variables $u=(u_j)_{j=1}^n, v=(v_j)_{j=1}^n$, define the Boolean polynomial $\mathrm{PairwiseDistinct}(u,v)$
$$
\mathrm{PairwiseDistinct}(u,v)=1+\prod_{j=1}^n (1+u_j+v_{j}).
$$

Then the sets $u|_S$ and $v|_S$ are distinct if and only if 
$$\mathrm{PairwiseDistinct}(u,v)|_S=1.$$

\end{definition} 

\begin{definition} 

Define $\mathrm{Distinct}(X)  \in B$ by
$$
\mathrm{Distinct}(X) =\prod_{1 \leq i_1<i_2 \leq m}\mathrm{PairwiseDistinct}(X_{i_1},X_{i_2}). 
$$
\end{definition} 

\begin{definition} 
A set $u|_S$ is a member of the set $X|_S$ if and only if 
$$
\prod_{i=1}^m \mathrm{PairwiseDistinct}(u,X_i)|_S = 0.
$$ 
Define $\mathrm{Member}(X,u)$ by 
$$
\mathrm{Member}(X,u) = 1+ \prod_{i=1}^m \mathrm{PairwiseDistinct}(u,X_i)
$$
\end{definition} 

\begin{definition} 

Define the sequence $u \cap v$ 
$$
u \cap v= (u_jv_j)_{j=1}^n
$$

Thus define the Boolean polynomial $\mathrm{IntersectionClosed}(X)$ by 
$$
\mathrm{IntersectionClosed}(X) = \prod_{1 \leq i_1<i_2 \leq m}\mathrm{Member}(X,X_{i_1} \cap X_{i_2}) 
$$

Therefore the set $X|_S$ is intersection-closed if and only if 
$$
\mathrm{IntersectionClosed}(X)|_S = 1.
$$
\end{definition} 

\begin{definition} 
Define a matching on the set $[1,2k]$ to be a set partition of $[1,2k]$ with each set having size 2. Define a matching on the set $[1,2k+1]$ to be a set partition of $[1,2k+1]$ with one set of size 1 and all other sets of size 2. Let $\mathrm{Matchings}(m)$ denote the set of all matchings of $[1,m]$. For $M \in  \mathrm{Matchings}(m)$, define
$$
H_{M,j}(X) = \begin{cases} \prod_{\{ i_1,i_2\} \in M} (1+ x_{i_1,j}x_{i_2,j}) &\text{ if $m$ is even} \\ 
(1+x_{i_0, j})\prod_{\{ i_1,i_2\} \in M} (1+ x_{i_1,j}x_{i_2,j}) &\text{ if $m$ is odd} 
\end{cases}
$$
where $\{i_0\}$ is the unique set of size 1 in $M$, if $m$ is odd. 

For a specialization $S$ of $X$ and integer $j, 1 \leq j \leq n$, define the statement $H(S,j)$ to be: 
\[
\#\{ i \colon j \in X_i|_S\} \leq  \lfloor \frac{m}{2}\rfloor
\]
 where $m$ is the number of sets in the set $X|_S$. 
\end{definition}

\begin{remark}If statement H$(S,j)$ is true, then there exists some matching $M$ of $[1,m]$ such that  
$$
H_{M,j}(X)|_S = 1.
$$

Thus if statement $H(S,j)$ is true, then 
$$
\prod_{M \in \mathrm{Matchings}(m)} (1+H_{M,j}(X))|_S=0. 
$$
\end{remark}

\begin{definition} 
Define $$\mathrm{HalfMembership}(X) = \prod_{j=1}^n \prod_{M \in \mathrm{Matchings}(m)} (1+H_{M,j}(X)).$$
\end{definition}

\begin{definition} 
Let $\mathrm{Specializations}_{\mathrm{ICC}}(m,n)$ denote the set of specializations $S$ of the variables $x_{i,j}$ such that the set $X|_S$ is an intersection-closed set of distinct sets. Define the statement $\mathrm{ICC}(m,n)$ to be ``For a specialization $S \in \mathrm{Specializations}_{\mathrm{ICC}}(m,n)$, there exists some integer $j$ such that 
\[
\#\{ i \colon j \in X_j|_S \} \leq \lfloor \frac{m}{2}\rfloor.  
\]
\end{definition}


\begin{theorem} 
Statement $\mathrm{ICC}(m,n)$ is true if and only if $\mathrm{ICC}_{m,n}(X)$ is the zoo Boolean polynomial. 
\end{theorem}

\begin{proof}
If the statement ICC$(m,n)$ is true, then for any $S \in \mathrm{Specializations}_{\mathrm{ICC}}(m,n)$. 
we have
$$
\begin{aligned} 
\mathrm{Distinct}(X)|_S&=1\\ 
\mathrm{IntersectionClosed}(X)|_S&=1 \\
\mathrm{HalfMembership}(X)|_S&=0.
\end{aligned}
$$
And if $S \notin \mathrm{Specializations}_{\mathrm{ICC}}(m,n)$, then either 
$$
\mathrm{Distinct}(X)|_S=0 \text{ or } \mathrm{IntersectionClosed}(X)|_S=0. 
$$

So, if statement ICC$(m,n)$ is true, then the polynomial
$$ \label{7}
\mathrm{ICC}_{m,n}(X) = \mathrm{Distinct}(X)
\mathrm{IntersectionClosed}(X)
\mathrm{HalfMembership}(X)
$$
evaluates to 0 for all specializations $S$, and thus must be the zero Boolean polynomial. 

Conversely, suppose $\mathrm{ICC}_{m,n}(X)$ is the zero Boolean polynomial. Then for $S \in \mathrm{Specializations}_{\mathrm{ICC}}(m,n)$, we know 
$$ \mathrm{Distinct}(X)|_S = 
\mathrm{IntersectionClosed}(X)|_S = 1,
$$ so must have 
$$\mathrm{HalfMembership}(X)|_S=0.$$ Thus $H_{M,j}(X)|_S=1$ for some $M$ and $j$, which implies that the element $j$ is in at most $\lfloor \frac{m}{2} \rfloor$ sets of $X|_S$. Thus the statement ICC$(m,n)$ is true if and only if $\mathrm{ICC}_{m,n}(X)$ is the zero Boolean polynomial. This completes the proof.
\end{proof}

\section{Formulas for Member$(X,u)$ and  \newline Member$(X,u)$PairwiseDistinct$(u,v)$}


\end{document}
